\section{An unfinished attempt at QGA self-consistency}
\label{sec:self}
\label{sec:semantics}

This appendix records an attempt to prove \QGA's consistency within
\QGA{} itself.  The attempt was not completed and has been abandoned.
I retain this account in case its ideas or difficulties are useful to
future collaborators.  The experimental files
\texttt{pure/QGA\_on\_GA.thy} and \texttt{pure/qga\_encode.thy} remain in
the repository, but they contain unproved lemmas and do not establish
\QGA{} self-consistency.

\subsection{The proposed semantics}
\label{sec:semantics:formulas}
\label{sec:semantics:quant}

The aim was to repeat the \BGA{} construction with encoded \QGA{} formulas
and proofs.  A formula could evaluate to true, to false, or to neither;
the last case would represent an ungrounded formula.  Equality would compare
terminating terms, negation would exchange truth and falsity, and
disjunction would be true when either disjunct was true and false when both
were false.

Quantifiers were the main difficulty.  Reading a universal as true when
every numeral instance is true leaves infinitely many instances to check,
even with a fixed evaluation budget.  The alternative explored was to
search for a proof of the body for an arbitrary natural number, and for
a refutation of an instance to establish falsity.  This was motivated by
reflective grounded arithmetic~\cite{ford-rga}.  Bounding proof search
would permit finite checks, but relating their results to truth would
require further proofs: a certificate cannot be treated as sound merely
because the checker accepts it.

\subsection{Ideas retained from the attempt}

The \BGA{} proof suggested keeping hypothesis satisfaction decidable at a
fixed budget, while allowing a less restrictive predicate for conclusions.
The attempt explored an additional ``witnessed'' semantics, in which a
quantifier certificate was accompanied by the semantic facts it was meant
to establish.  Separate searches for truth and falsity were also considered,
to avoid making one search take priority over the other.

Another proposal allowed a quantifier certificate to retain closed
hypotheses, provided their satisfaction was checked within a recorded
budget.  The intention was to reuse a derivation already present in the
proof being checked, rather than first turn every satisfied hypothesis
into a proof from no hypotheses.

These were proposed ways to organize the argument, not established
solutions.  In particular, the necessary substitution and soundness
arguments were not completed.
